\scanpage{5}
\margin{. . . .}
\begin{center}- 180 -\end{center}
\hfill n° 391

Alors, (§ 1, n° 3, prop. 8), on aurait $\sum_{k\ne i}q_{ik}^2\geq\frac14+(\frac{\sqrt3}{2})^2=1$, ce qui contredit le lemme 1.

COROLLAIRE 4.-- \uline{Un sommet ne peut appartenir à deux arêtes distinctes d'ordre} $\geq4$.

En effet, (§ 1, n° 3, prop. 8), on a $(\frac{\sqrt2}{2})^2+(\frac{\sqrt2}{2})^2=1$.

Soit $\{i,j\}$ une arête de $X$ ; soit $\underline{M}'=(m'_{hk})$, $(1\leq h\leq\ell,\ 1\leq k\leq\ell,\ h\ne j,\ k\ne j)$, la matrice d'ordre $\ell-1$ définie par $m'_{hk}=m_{hk}$ si $h\ne i$ ou $k\ne i$, $m'_{ik}=m'_{ki}=m_{ik}+m_{jk}$ pour tout $k$ distinct de $i$ et de $j$, $m'_{ii}=m_{ii}=1$. On dit que le schéma de Coxeter $X'$ de $\underline M'$ s'obtient à partir de $X$ \uline{identifiant} $j$ \uline{à} $i$. Les arêtes de $X'$ sont donc les ensembles $\{i,k\}$, $(k\ne j)$, qui sont des arêtes de $X$, les ensembles $\{i,k'\}$, $(k'\ne j)$, tels que $\{j,k'\}$ soit une arête de $X$, et les arêtes $\{k,k'\}$ de $X$ liant des sommets distincts de $i$ et de $j$. Si $k$ est distinct de $i$ et de $j$, on notera que $\{i,k\}$ et $\{j,k\}$ ne peuvent tous deux être des arêtes de $X$, sans quoi elles formeraient avec $\{i,j\}$ un circuit dans $X$ ; on en conclut que si $k$ est \uline{point de ramification} dans $X$, il l'est encore dans $X'$.

LEMME 2.-- \uline{Si l'arête} $\{i,j\}$ \uline{est d'ordre} 3, $G(\underline M')$ \uline{est un groupe fini}.

Notons en effet que si $k$ est distinct de $i$ et de $j$, on ne peut avoir \uline{à la fois} $m_{ik}\geq3$ et $m_{jk}\geq3$, sans quoi les arêtes $\{i,j\}$, $\{j,k\}$, $\{k,i\}$ formeraient un circuit dans $X$. Si on pose $q'_{hk}=-\cos(\pi/m'_{hk})$, $(h,k\text{ distincts de }j)$, on a donc $q'_{ik}=q_{ik}+q_{jk}$ pour $k$ distinct de $i$ et de $j$ ; comme $q_{ii}=q_{jj}=1$ et $q_{ij}=-\frac12$ par hypothèse, on voit que la forme quadratique $B_{\underline M'}$ est telle que pour $x=\sum_{k\ne j}\xi_k e_k$, $B_{\underline M'}(x,x)$ soit égal au nombre obtenu en remplaçant $\xi_j$ par $\xi_i$ dans $q(\xi_1,\ldots,\xi_\ell)$ ; on a donc $B_{\underline M'}(x,x)>0$ par hypothèse si $x\ne0$.
\mspage{4}
